% DeepPass21: continuum normalization from a finite box.

把无限空间暂时放进边长为 $L$ 的周期立方盒，可以把连续动量谱变成离散格点。周期边界条件要求
\begin{equation}
 \boxed{
 p_i=\frac{2\pi\hbar}{L}n_i,
 \qquad n_i\in\mathbb Z,
 \qquad
 \Delta p_i=\frac{2\pi\hbar}{L}.}
 \label{eq:periodic-box-momenta-dp68}
\end{equation}
因此每个允许动量点在三维动量空间中占据体积 $(2\pi\hbar/L)^3$。当 $L\to\infty$ 时，离散 Riemann 和与位置完备核分别趋于
\begin{equation}
\boxed{
\sum_{\mathbf n}\longrightarrow
\frac{L^3}{(2\pi\hbar)^3}\int\dd^3 p,
\qquad
\frac1{L^3}\sum_{\mathbf n}\ee^{\ii\mathbf p_{\mathbf n}\cdot(\mathbf r-\mathbf r')/\hbar}
\longrightarrow\delta^{(3)}(\mathbf r-\mathbf r').}
\label{eq:box-to-continuum-dp21}
\end{equation}
这里第二个极限按分布意义理解；它不是逐点函数极限。

\begin{derivation}[从\eqnrefs{eq:periodic-box-momenta-dp68}到\eqnrefs{eq:box-to-continuum-dp21}]
对随动量变化足够平滑的测试函数 $F(\mathbf p)$，

\begin{align*}
 \sum_{\mathbf n}F(\mathbf p_{\mathbf n})
 &=\sum_{\mathbf n}
 \left(\frac{2\pi\hbar}{L}\right)^3
 \frac{F(\mathbf p_{\mathbf n})}{(2\pi\hbar/L)^3}\\
 &\longrightarrow
 \frac{L^3}{(2\pi\hbar)^3}
 \int\dd^3 p\,F(\mathbf p).
\end{align*}
这给出第一条极限。

盒归一化平面波为
\begin{equation*}
 \langle\mathbf r|\mathbf n\rangle
 =L^{-3/2}\ee^{\ii\mathbf p_{\mathbf n}\cdot\mathbf r/\hbar},
\end{equation*}
离散完备性
$\sum_{\mathbf n}|\mathbf n\rangle\langle\mathbf n|=I$
在位置表象中变成
\begin{equation*}
 \frac1{L^3}\sum_{\mathbf n}
 \ee^{\ii\mathbf p_{\mathbf n}\cdot(\mathbf r-\mathbf r')/\hbar}
 =\langle\mathbf r|I|\mathbf r'\rangle.
\end{equation*}
将刚得到的求和--积分极限用于左端，
\begin{equation*}
 \frac1{L^3}\sum_{\mathbf n}
 \ee^{\ii\mathbf p_{\mathbf n}\cdot\Delta\mathbf r/\hbar}
 \longrightarrow
 \int\frac{\dd^3 p}{(2\pi\hbar)^3}
 \ee^{\ii\mathbf p\cdot\Delta\mathbf r/\hbar}
 =\delta^{(3)}(\Delta\mathbf r),
\end{equation*}
即正文\eqnrefs{eq:box-to-continuum-dp21} 的第二条关系。
\end{derivation}


